\documentclass[11pt,a4paper]{article}
% XeLaTeX-friendly preamble
\usepackage{fontspec}
\setmainfont{TeX Gyre Termes}
\usepackage[brazil]{babel}
\usepackage{geometry}
\geometry{margin=2.5cm}
\usepackage{amsmath,amssymb}
\usepackage{enumitem}
\usepackage{hyperref}
\usepackage{microtype}
\title{Prova de Revis\~ao — Engenharia de Reservat\'orios I}
\author{}
\date{}
\begin{document}
\maketitle
\thispagestyle{empty}
\noindent\textbf{Instru\c{c}\~oes:} Responda a todas as quest\~oes. Use os espa\c{c}os fornecidos para rascunho e c\'alculos.
\vspace{1em}

\section*{1. Verdadeiro ou falso — Cap\'itulo 1 (V/F + justifique em at\'e 2 linhas)}
\begin{enumerate}[label=\alph*)]
\item Em reservat\'orios leves, a viscosidade do \'oleo costuma ser menor que em \'oleos pesados. \hfill \textbf{(V/F)}\\
\vspace{2\baselineskip}
\item A presen\c{c}a de g\'as dissolvido no \'oleo sempre aumenta a mobilidade do fluido. \hfill \textbf{(V/F)}\\
\vspace{2\baselineskip}
\item Uma drive de \'agua (water drive) mant\'em a press\~ao do reservat\'orio por influxo de \'agua. \hfill \textbf{(V/F)}\\
\vspace{2\baselineskip}
\item A saturac\~ao relativa da \'agua \'e igual a 1 menos a saturac\~ao do \'oleo. \hfill \textbf{(V/F)}\\
\vspace{2\baselineskip}
\end{enumerate}

\section*{2. Disserta\c{c}\~ao (Cap\'itulo 4)}
Tema: ``Incertezas no m\'etodo volum\'etrico: fontes, quantifica\c{c}\~ao e mitiga\c{c}\~ao''.

Instru\c{c}\~oes: Elabore uma reda\c{c}\~ao t\'ecnica organizada em Introdu\c{c}\~ao, Desenvolvimento e Conclus\~ao. Discuta pelo menos duas fontes de incerteza e descreva um m\'etodo pr\'atico para quantific\'a-las (m\'ax. 1 p\'agina).

\vspace{8\baselineskip}

\section*{3. C\'alculo/pr\'atico — Cap\'itulo 2}
Enunciado: Um ensaio PVT fornece: $V_{res}=1.20$ bbl; $V_{surf}=1.00$ STB; $R_s=400$ scf/STB.

a) Calcule o fator de forma\c{c}\~ao do \'oleo $B_o$ (bbl/STB).\\
b) Converta $R_s$ para m$^3$/m$^3$ (use $1\ \mathrm{scf}=0.0283168\ \mathrm{m}^3$, $1\ \mathrm{STB}=0.1589873\ \mathrm{m}^3$).

\vspace{8\baselineskip}

\section*{4. C\'alculo/pr\'atico — Cap\'itulo 3}
Enunciado: Amostra saturada: $m_{sat}=130\ \mathrm{g}$; $m_{dry}=105\ \mathrm{g}$; $\rho_{oil}=0.84\ \mathrm{g/cm}^3$; $V_{tot}=180\ \mathrm{cm}^3$.

a) Calcule o volume de fluido nos poros $V_f$ e a porosidade $\phi$ (\%).

\vspace{8\baselineskip}

\bigskip
\noindent Boa sorte.

\end{document}
\documentclass[11pt]{article}
\usepackage{fontspec}
\setmainfont{TeX Gyre Termes}
\usepackage[brazilian]{babel}
\usepackage{geometry}
\geometry{a4paper,margin=25mm}
\usepackage{enumitem}
\usepackage{amsmath,amssymb}
\usepackage{setspace}
\onehalfspacing
\title{Prova — Capítulos 1–4}
\date{}
\begin{document}
\maketitle
\noindent \textbf{Instruções}: Responda as quatro tarefas abaixo. Mostre todos os cálculos e justifique as respostas.

\begin{enumerate}
\item \textbf{Verdadeiro ou falso — Capítulo 1.} Para cada afirmação, marque V ou F e justifique em até duas linhas.
\begin{enumerate}[label=\alph*)]
\item Em reservatórios leves, o óleo costuma ter baixa densidade e alta mobilidade.
\vspace{2\baselineskip}
\item A presença de gás dissolvido no óleo sempre aumenta a mobilidade do fluido.
\vspace{2\baselineskip}
\item Um \emph{water drive} mantém pressão por influxo do aquífero.
\vspace{2\baselineskip}
\item A condensação de hidrocarbonetos no reservatório ocorre se a pressão cair abaixo do ponto de orvalho.
\vspace{2\baselineskip}
\item A saturação relativa da água é igual a $1$ menos a saturação do óleo.
\vspace{2\baselineskip}
\end{enumerate}

\newpage
\item \textbf{Dissertação — Capítulo 4.} Tema: \emph{Incertezas no método volumétrico: fontes, quantificação e mitigação}. Escreva uma redação técnica (máx. 1 página) estruturada em Introdução, Desenvolvimento e Conclusão.
\vspace{14\baselineskip}

\newpage
\item \textbf{Cálculo/prático — Capítulo 2 (PVT).} Enunciado: Um ensaio PVT fornece: $V_{res}=1.20$ bbl; $V_{surf}=1.00$ STB; $R_s=400$ scf/STB. \\
(a) Calcule o fator de formação do óleo $B_o$ em bbl/STB. \\
(b) Converta $R_s$ para m$^3$/m$^3$ (use $1\ \text{scf}=0.0283168\ \text{m}^3$, $1\ \text{STB}=0.1589873\ \text{m}^3$). \\
Espaço para cálculos:
\vspace{12\baselineskip}

\newpage
\item \textbf{Cálculo/prático — Capítulo 3 (Rochas).} Enunciado: Amostra saturada: $m_{sat}=130\ \text{g}$; $m_{dry}=105\ \text{g}$; $\rho_{oil}=0.84\ \text{g/cm}^3$; $V_{tot}=180\ \text{cm}^3$. \\
(a) Calcule o volume de fluido nos poros $V_f$ e a porosidade $\phi$ (\%). \\
(b) Apresente os passos e respostas com 3 algarismos significativos. \\
Espaço para cálculos:
\vspace{12\baselineskip}

\end{enumerate}

\end{document}
